\documentclass[a4paper,14pt,oneside,final]{extreport}
\usepackage[utf8]{inputenc}
\usepackage[T2A]{fontenc}
\usepackage[russian]{babel}
\usepackage{graphicx}
\graphicspath{{png/}}
\usepackage{geometry}
\usepackage{fancyhdr}
\usepackage{tabularx}
\usepackage{amsmath}
\usepackage{amsfonts}
\usepackage{amssymb}
\usepackage{array}
\usepackage{longtable}
\usepackage{url}
\usepackage{textcomp}
\usepackage{indentfirst}
\usepackage{enumitem}
\usepackage{listings}
\usepackage{xcolor}
\usepackage{caption}
\usepackage{float}
\usepackage{tocloft}

\geometry{left=3cm,top=2.0cm,right=1.5cm,bottom=2.7cm}
\sloppy
\frenchspacing
\raggedbottom
\setlength{\parindent}{1.25cm}

% Тире нормальной «тонкости» (не жирная плашка)
\newcommand{\mdash}{%
	\unskip\nobreak\hskip0.2em\relax
	{\raisebox{0.45ex}{\rule{0.7em}{0.35pt}}}%
	\nobreak\hskip0.2em\relax
}

\lstset{
	language={},
	basicstyle=\small\ttfamily,
	keywordstyle=\bfseries,
	commentstyle=\itshape,
	numbers=left,
	numberstyle=\tiny,
	stepnumber=1,
	numbersep=5pt,
	showspaces=false,
	showstringspaces=false,
	frame=single,
	breaklines=true,
	breakatwhitespace=true,
	tabsize=4,
	columns=fullflexible,
	keepspaces=true,
	captionpos=b
}

\renewcommand{\contentsname}{Содержание}
% Заголовок TOC как у <<Перечень...>>: sectioncentered*; сам tocloft-title отключаем
\makeatletter
\renewcommand{\@cftmaketoctitle}{}
\makeatother
% Ровная табуляция номеров в содержании
\newlength{\pziistocnum}
\setlength{\pziistocnum}{2.6em}
\cftsetindents{section}{0em}{\pziistocnum}
\cftsetindents{subsection}{\pziistocnum}{\pziistocnum}
\cftsetindents{subsubsection}{2\pziistocnum}{\pziistocnum}
\renewcommand{\cftsecfont}{\normalfont}
\renewcommand{\cftsubsecfont}{\normalfont}
\renewcommand{\cftsubsubsecfont}{\normalfont}
\renewcommand{\cftsecpagefont}{\normalfont}
\renewcommand{\cftsubsecpagefont}{\normalfont}
\renewcommand{\cftsubsubsecpagefont}{\normalfont}

\renewcommand{\lstlistingname}{Листинг}
\addto\captionsrussian{\renewcommand{\lstlistingname}{Листинг}}

\DeclareCaptionLabelFormat{stbfigure}{Рисунок #2}
\DeclareCaptionLabelFormat{stbtable}{Таблица #2}
\DeclareCaptionLabelFormat{stblisting}{Листинг #2}
\DeclareCaptionLabelSeparator{stb}{~\mdash~}
% Подписи снизу, курсивом, по центру
\captionsetup{
	labelsep=stb,
	font={small,it},
	labelfont={small,it},
	justification=centering,
	singlelinecheck=true,
	position=bottom
}
\captionsetup[figure]{labelformat=stbfigure}
\captionsetup[lstlisting]{labelformat=stblisting}
\captionsetup[table]{labelformat=stbtable}

\makeatletter
\renewcommand{\thesection}{\arabic{section}}
\makeatother
\setcounter{secnumdepth}{3}
\setcounter{tocdepth}{3}

\AtBeginDocument{\numberwithin{equation}{section}}
\AtBeginDocument{\numberwithin{table}{section}}
\AtBeginDocument{\numberwithin{figure}{section}}
\AtBeginDocument{\numberwithin{lstlisting}{section}}

\pagestyle{fancy}
\fancyhf{}
\fancyfoot[R]{\thepage}
\renewcommand{\footrulewidth}{0pt}
\renewcommand{\headrulewidth}{0pt}
\fancypagestyle{plain}{\fancyhf{}\rfoot{\thepage}}

\makeatletter
\renewcommand\section{%
	\clearpage\@startsection {section}{1}%
	{1.25cm}%
	{-1em \@plus -1ex \@minus -.2ex}%
	{1em \@plus .2ex}%
	{\raggedright\hyphenpenalty=10000\normalfont\large\bfseries\MakeUppercase}}
\renewcommand\subsection{%
	\@startsection{subsection}{2}%
	{1.25cm}%
	{-1em \@plus -1ex \@minus -.2ex}%
	{1em \@plus .2ex}%
	{\raggedright\hyphenpenalty=10000\normalfont\normalsize\bfseries}}
\renewcommand\subsubsection{%
	\@startsection{subsubsection}{3}%
	{1.25cm}%
	{-1em \@plus -1ex \@minus -.2ex}%
	{1em \@plus .2ex}%
	{\raggedright\hyphenpenalty=10000\normalfont\normalsize\bfseries}}
\newcommand\sectioncentered{%
	\clearpage\@startsection {section}{1}%
	{\z@}%
	{-1em \@plus -1ex \@minus -.2ex}%
	{1em \@plus .2ex}%
	{\centering\hyphenpenalty=10000\normalfont\large\bfseries\MakeUppercase}}
\makeatother

\setlist{nosep,topsep=0.35em,partopsep=0pt,parsep=0pt}
\makeatletter
\AddEnumerateCounter{\asbuk}{\@asbuk}{щ)}
\makeatother
\renewcommand{\labelenumi}{\asbuk{enumi})}
\renewcommand{\labelenumii}{\arabic{enumii})}
% Списки без красной строки: у левого края текста, маркер --- дефис
\setlist[itemize]{%
	leftmargin=1.1em,%
	itemindent=0pt,%
	labelsep=0.35em,%
	label={\normalfont{-}}%
}
\setlist[enumerate,1]{%
	leftmargin=1.4em,%
	itemindent=0pt,%
	labelsep=0.35em%
}
\setlist[enumerate,2]{%
	leftmargin=2.6em,%
	itemindent=0pt%
}

\makeatletter
\renewenvironment{thebibliography}[1]
{\sectioncentered*{Список использованных источников}
	\@mkboth{\MakeUppercase\refname}{\MakeUppercase\refname}%
	\list{\@biblabel{\@arabic\c@enumiv}}%
	{\settowidth\labelwidth{\@biblabel{#1}}%
		\setlength{\itemindent}{\dimexpr\labelwidth+\labelsep+1em}
		\leftmargin\z@
		\@openbib@code
		\usecounter{enumiv}%
		\let\p@enumiv\@empty
		\renewcommand\theenumiv{\@arabic\c@enumiv}}%
	\sloppy
	\clubpenalty4000
	\@clubpenalty \clubpenalty
	\widowpenalty4000%
	\sfcode`\.\@m}
{\def\@noitemerr
	{\@latex@warning{Empty `thebibliography' environment}}%
	\endlist}
\makeatother
